\section{SELECT Dasar dan WHERE}

Perintah \code{SELECT} adalah perintah SQL yang paling sering digunakan --- berfungsi untuk mengambil (membaca) data dari satu atau lebih tabel. Sintaks dasar: \code{SELECT kolom1, kolom2 FROM tabel;} untuk memilih kolom tertentu, atau \code{SELECT * FROM tabel;} untuk mengambil semua kolom. Meskipun \code{SELECT *} praktis untuk eksplorasi, di aplikasi produksi sebaiknya disebutkan kolom secara eksplisit agar query lebih efisien dan tidak mengambil data yang tidak diperlukan \cite{mysqlDocs}.

Klausa \code{WHERE} memfilter baris yang dikembalikan berdasarkan kondisi. Kondisi dapat menggunakan operator perbandingan (\code{=}, \code{<>}, \code{<}, \code{>}, \code{<=}, \code{>=}), operator logika (\code{AND}, \code{OR}, \code{NOT}), serta operator khusus seperti \code{BETWEEN}, \code{IN}, \code{LIKE}, dan \code{IS NULL}. Operator \code{LIKE} mendukung wildcard: \code{\%} untuk nol atau lebih karakter, \code{\_} untuk satu karakter. Kombinasi \code{WHERE} dengan operator ini memungkinkan pencarian data yang sangat fleksibel \cite{mysqlGettingStarted}.

\begin{webcode}[language=SQL, caption={Contoh SELECT dengan WHERE}]
-- Semua produk
SELECT * FROM produk;

-- Kolom tertentu
SELECT nama, harga FROM produk;

-- Filter dengan WHERE
SELECT nama, harga FROM produk WHERE harga > 100000;

-- Kombinasi kondisi
SELECT nama, harga, stok FROM produk
  WHERE kategori = 'elektronik' AND stok > 0;

-- Pencarian teks dengan LIKE
SELECT * FROM produk WHERE nama LIKE '%laptop%';

-- Nilai dalam daftar
SELECT * FROM produk WHERE kategori IN ('elektronik', 'pakaian');

-- Rentang nilai
SELECT * FROM produk WHERE harga BETWEEN 50000 AND 200000;
\end{webcode}

